% ==============================================================================
% Plantilla Institucional de Trabajo Terminal — UPIIZ IPN
% Archivo: capitulos/03-antecedentes.tex
% ==============================================================================
% LÍMITE NORMATIVO INSTITUCIONAL: 1 cuartilla (aproximadamente 1 página).
% ==============================================================================

\chapter{Antecedentes}
\label{cap:antecedentes}

En este capítulo se expone el contexto histórico y la evolución tecnológica previa del problema de estudio, identificando los antecedentes científicos directos y los proyectos desarrollados previamente a nivel internacional, nacional y dentro del Instituto Politécnico Nacional.

Se sugiere redactar este capítulo considerando las siguientes pautas:
\begin{itemize}
    \item \textbf{Evolución tecnológica del área:} Describa cómo han evolucionado históricamente los métodos, dispositivos y sistemas empleados para resolver la problemática planteada.
    \item \textbf{Trabajos previos en la UPIIZ-IPN:} Mencione si existen Trabajos Terminales, proyectos de investigación o bancos de laboratorio desarrollados previamente en la academia que sirvan como punto de partida o referencia directa para esta propuesta.
    \item \textbf{Limitaciones de los desarrollos previos:} Explique claramente qué aspectos no fueron cubiertos o qué áreas de oportunidad motivan el presente Trabajo Terminal (por ejemplo: falta de lazo cerrado, procesamiento en tiempo real limitado, requerimientos de portabilidad o costo excesivo).
\end{itemize}

\noindent La redacción debe mantenerse en pasado para investigaciones y proyectos concluidos, y en presente para el estado actual de los hechos.
